 Στην παρούσα διπλωματική εργασία αναπτύσσεται και αξιολογείται μια μεθοδολογία βασισμένη σε δεδομένα (\textenglish{data-driven}) για την εκτίμηση και τη συσταδοποίηση των ρυθμών ταλάντωσης Συστημάτων Ηλεκτρικής Ενέργειας (ΣΗΕ), αξιοποιώντας δεδομένα κανονικής λειτουργίας (\textenglish{ambient data}). Για την εκτίμηση των συχνοτήτων και των συντελεστών απόσβεσης των ρυθμών εφαρμόζεται η μέθοδος αναγνώρισης υποχώρου \textenglish{N4SID}, ενώ αξιολογούνται επτά αλγόριθμοι συσταδοποίησης ως προς την ικανότητά τους να ομαδοποιούν τις εκτιμήσεις και να αναγνωρίζουν τους κυρίαρχους ρυθμούς ταλάντωσης του συστήματος.

Η προτεινόμενη μεθοδολογία εφαρμόζεται στο πρότυπο σύστημα IEEE 39-Bus, ενώ συμπληρωματικά πραγματοποιείται ανάλυση της απόκρισης σε διαταραχές με τη μέθοδο \textenglish{Matrix Pencil} στα πρότυπα συστήματα \textenglish{Kundur Two-Area} και \textenglish{IEEE 39-Bus}. Τα αποτελέσματα της συγκριτικής αξιολόγησης αναδεικνύουν τον αλγόριθμο \textenglish{DBSCAN} ως τον αποτελεσματικότερο ως προς τη συμφωνία των εκτιμώμενων ρυθμών με τους αντίστοιχους ρυθμούς αναφοράς του λογισμικού \textenglish{PowerFactory}. Αντίστοιχα, ο \textenglish{Hierarchical DBSCAN (HDBSCAN)} παρουσιάζει τον καλύτερο συμβιβασμό μεταξύ ποιότητας συσταδοποίησης, κάλυψης των δεδομένων και απόρριψης θορύβου.

Τέλος, διερευνάται η επίδραση της διείσδυσης αιολικών μονάδων με επαγωγικές γεννήτριες διπλής τροφοδότησης (\textenglish{DFIG}) στους ρυθμούς ταλάντωσης και στην αποτελεσματικότητα των εξεταζόμενων αλγορίθμων. Τα αποτελέσματα καταδεικνύουν μεταβολές στις συχνότητες των ρυθμών και στη σχετική συμμετοχή των σύγχρονων γεννητριών, χωρίς ωστόσο να παρατηρούνται ουσιαστικές διαφοροποιήσεις στη συγκριτική απόδοση των μεθόδων συσταδοποίησης

\paragraph{Λέξεις-κλειδιά:} ανάλυση διαταραχών, ανάλυση κανονικής λειτουργίας, διασυνδετικοί ρυθμοί, ρυθμοί ταλάντωσης, συσταδοποίηση, ταυτοποίηση συστημάτων, τοπικοί ρυθμοί, \textenglish{Matrix Pencil}, N4SID
